\section{Introduction}\label{sec:introduction}

Let $(M^n,g)$ be a smooth closed connected Riemannian manifold. We use
$\Delta_g=\diver_g\nabla_g$ and write
\[
 -\Delta_g u=\lambda u,\qquad
 \Z(u)=\{x\in M:u(x)=0\}.
\]
Yau posed the nodal-size question in his problem collection
\cite{Yau1982}. In its smooth closed-manifold formulation, Yau's
nodal-set conjecture predicts that the $(n-1)$-dimensional
Hausdorff measure of $\Z(u)$ is bounded above and below by positive
constant multiples of $\sqrt\lambda$, with the constants depending only
on the fixed manifold and metric \cite[p.~162]{DonnellyFefferman1988}.
The upper assertion is required to
hold for every nonzero real eigenfunction, including all eigenfunctions
in a multiple eigenspace.

We disprove that upper assertion in the smooth category. In fact, the
assertion that one has an upper bound with every positive power loss,
\begin{equation}\label{eq:near-sharp-target}
 \Hh_g^{n-1}(\Z(u))
 \le C(M,g,\epsilon)\lambda^{1/2+\epsilon}
 \qquad\text{for every }\epsilon>0,
\end{equation}
fails.

\begin{theorem}\label{thm:main}
There exist a smooth Riemannian metric $g$ on $S^4\times S^1$, a number
$\epsilon_0>0$, and nonzero real smooth functions $u_j$ with positive
eigenvalues $\lambda_j$ such that
\[
 -\Delta_g u_j=\lambda_j u_j,\qquad
 \lambda_j\longrightarrow\infty,\qquad
 \frac{\Hh_g^4(\Z(u_j))}
      {\lambda_j^{1/2+\epsilon_0}}\longrightarrow\infty.
\]
The metric and $\epsilon_0$ are fixed throughout the sequence.
\end{theorem}

In particular, \eqref{eq:near-sharp-target} fails for every
$0<\epsilon\le\epsilon_0$ for the metric of Theorem~\ref{thm:main}.

The theorem concerns ordinary Laplace--Beltrami eigenfunctions on one
closed connected manifold. Taking a product with a fixed closed manifold
gives the same conclusion in each dimension greater than five. The
construction makes no assertion about the corresponding upper bound in
dimensions three or four.

In dimensions three and four, a companion paper constructs fixed smooth
metrics and real eigenfunction sequences with $\lambda\to\infty$ for which
the nodal measure divided by $\sqrt\lambda$ tends to infinity, without
asserting a fixed positive power-law excess
\cite[Theorems~1.1 and~1.2]{OpenAINodalQualitative2026}.

\paragraph{Background.}
Donnelly and Fefferman proved the sharp two-sided nodal estimate for
real-analytic metrics \cite[Theorem~1.2]{DonnellyFefferman1988}.
For smooth metrics, Logunov proved the sharp lower bound
\cite[Theorem~1.1]{LogunovLower2018} and, in dimensions at least three,
a polynomial upper bound with a dimension-dependent exponent
\cite[Theorem~6.2]{LogunovUpper2018}.
Regularity assumptions between analyticity and smoothness also yield
stronger upper estimates: Hezari proved a $\lambda^{s/2}$ bound in
Gevrey class $s\ge1$, and a $\sqrt\lambda\log(\lambda+e)$ bound for a
specified quasianalytic class
\cite{Hezari2023}.
Those hypotheses do not cover arbitrary smooth metrics.
Theorem~\ref{thm:main} disproves the smooth upper-bound part of Yau's
conjecture and even the relaxed estimate \eqref{eq:near-sharp-target};
the lower-bound results remain compatible with it.

There is a complementary tradition of prescribing nodal geometry.
In dimensions at least three, Enciso and Peralta-Salas realize a smooth
closed connected orientable separating hypersurface as the nodal set of
a first nonconstant eigenfunction by choosing a metric
\cite{EncisoPeraltaSalas2015}. At high energy, Enciso, Peralta-Salas,
and Torres de Lizaur realize prescribed stable nodal components, up to
ambient diffeomorphism, on round spheres and flat tori, with a
localization restriction in the torus case
\cite{EncisoPeraltaSalasTorres2018}. Their construction already gives
arbitrarily high frequencies on fixed analytic metrics. The distinction
here is power-law excess in total nodal measure, which is not
furnished by those prescription theorems.

\paragraph{The quantitative mechanism.}
We first work with a symmetric elliptic operator on the four-sphere
whose positive principal tensor and positive density can be varied
independently. A local corrugation increases the integral of the
gradient of a logarithmic envelope. More specifically, the part of a
parent cube prepared for another corrugation has more new gradient
integral than the entire old parent cube. Repetition at a fixed scale
ratio therefore gives a multiplicative gain. The derivatives of order
$r$ grow at most exponentially in $(r+1)$ times the number of
generations. Let $k$ denote the construction frequency, with target
eigenvalue $k^2$. Taking a fixed small multiple of $\log k$ generations
produces a smooth real envelope $h=h_k$ on the base and a fixed
coordinate region $V$ such that
\[
 \int_V\sqrt{1+|\nabla h|^2}\,\dd x\ge c k^\gamma
\]
for one fixed $\gamma>0$, while keeping every fixed derivative of $h$
within a controlled power of $k$.

Finite Taylor solutions of the complex eikonal and transport equations
then give localized packets beneath $e^{kh}$. They achieve any prescribed
finite residual accuracy while preserving fixed derivative exponents.
A real Gaussian combination has a quantitatively nonvanishing first
jet everywhere and opposite signs on many pairs of small disks.
The first property permits an explicit simultaneous correction of the
principal tensor and density to make the function an exact
eigenfunction. The second gives stable lower bounds for nodal measure
without a regularity assumption on the zero set.

The distinction between constants and exponents is essential here.
At a given stage, constants and frequency thresholds may depend
arbitrarily badly on the required finite coefficient accuracy.
The exponent $\gamma$ does not; nor do the chart or coefficient
neighborhood. The positive constant in the nodal lower bound and the frequency threshold
may depend on the fixed input pair, collar, finite norm, and tolerance.
Consequently the construction can
meet every finite collection of earlier spectral-persistence tolerances
and still retain a fixed power gain. A summable iteration gives one
smooth limiting pair. Adding a warped circle realizes its weighted
operator as an ordinary Laplace--Beltrami operator.

These steps also give reusable local facts: a corrugation with a
repeatable integral gain, a phase-Hessian construction valid at zero
real slope, a polynomial first-jet noncancellation argument, and an
exact tensor--density correction. Each is proved below. In particular,
the argument does not appeal to a qualitative existence theorem with an
uncontrolled frequency threshold.

\paragraph{Organization.}
Section~\ref{sec:framework} states the quantitative one-step construction
and fixes the geometric notation.
Section~\ref{sec:envelopes} proves the envelope estimates.
Section~\ref{sec:waves} constructs and estimates the local packets.
Section~\ref{sec:global} produces exact eigenfunctions, proves the
one-step construction, and passes to the fixed metric in
Theorem~\ref{thm:main}.
